\section{All-order Pell--Lucas staircases and a paired companion splitter}
\label{sec:pell-lucas-all-order-staircase}

The two balancing-Pell channels admit a useful norm-one splice.  Put
\[
 (1+\sqrt2)^n=A_n+B_n\sqrt2,
 \qquad \alpha=3+2\sqrt2,\quad\beta=3-2\sqrt2,
\]
and let
\[
 u_0=0,\quad u_1=1,\quad u_{n+2}=6u_{n+1}-u_n,
\]
with companion $v_0=2,v_1=6,v_{n+2}=6v_{n+1}-v_n$.  Binet's
formulas give
\begin{equation}\label{eq:pell-lucas-splice}
 u_n=\frac{\alpha^n-\beta^n}{\alpha-\beta}
     =\frac{B_{2n}}2=A_nB_n,
 \qquad v_n=2A_{2n},
 \qquad v_n^2-32u_n^2=4.
\end{equation}
Thus the norm-one Lucas sequence sees the coprime channel product rather
than either channel alone.

We use Zhang's all-orders multiplication formula
\cite[Proposition~5.1]{ZhangLucas2026}.  This source is an arXiv v1
preprint, so we record the specialization and the additional argument in
full.  Fix an odd prime $\ell$, set
\[
 A=A_\ell,\qquad B=B_\ell,\qquad U=u_\ell=AB,
 \qquad \theta=\frac{\ell-1}{2},
\]
and define
\begin{equation}\label{eq:pell-lucas-coefficients}
 c_r(\ell)=
 \frac{\ell\prod_{j=1}^{r}
       (\ell^2-(2j-1)^2)}{4^r(2r+1)!},
 \qquad a_r=32^r c_r(\ell)\qquad(0\le r\le\theta).
\end{equation}
The source proves $c_r(\ell)\in\mathbf Z_{>0}$ and the exact identity
\begin{equation}\label{eq:pell-lucas-all-order-polynomial}
 Q_\ell:=\frac{u_{\ell^2}}{u_\ell}
       =\sum_{r=0}^{\theta}a_rU^{2r}.
\end{equation}

The prime-index rank congruences from the preceding Pell sections imply
that every prime $p\mid U$ satisfies
\begin{equation}\label{eq:pell-lucas-support-lower}
 p\equiv\left(\frac2p\right)\pmod{2\ell},
 \qquad p\ge2\ell-1.
\end{equation}
This lower bound makes every coefficient in
\eqref{eq:pell-lucas-all-order-polynomial} a unit on the actual support.

\begin{lemma}[All multiplication coefficients are support units]
\label{lem:pell-lucas-coefficient-unit}
For $0\le r\le\theta$,
\[
                         \gcd(a_r,U)=1.
\]
Moreover $c_\theta(\ell)=1$.
\end{lemma}

\begin{proof}
Let $p\mid U$.  For $1\le j\le r$, equation
\eqref{eq:pell-lucas-support-lower} gives
\[
 0<\ell-(2j-1)<p,
 \qquad 0<\ell+(2j-1)\le2\ell-2<p.
\]
Consequently $p$ divides none of the numerator factors in
\eqref{eq:pell-lucas-coefficients}.  Since the numerator equals the
denominator times the integer $c_r(\ell)$, the prime cannot divide
$c_r(\ell)$.  It also cannot divide $32$.  This proves the coprimality.

For the final coefficient write $\ell=2\theta+1$.  Then
\[
 \prod_{j=1}^{\theta}
    (\ell^2-(2j-1)^2)=4^\theta(2\theta)!,
\]
so the numerator and denominator of $c_\theta$ coincide.
\end{proof}

For $0\le r\le\theta$, subtract the first $r$ terms and put
\[
 D_r=Q_\ell-\sum_{i<r}a_iU^{2i},
 \qquad
 E_r=\sum_{i=r}^{\theta}a_iU^{2(i-r)}.
\]

\begin{theorem}[All-order support-unit staircase]
\label{thm:pell-lucas-all-order-staircase}
For every $0\le r\le\theta$,
\begin{equation}\label{eq:pell-lucas-tail-factorization}
 D_r=U^{2r}E_r,\qquad E_r\equiv a_r\pmod{U^2},
 \qquad \gcd(E_r,U)=1.
\end{equation}
Hence, for $p\mid U$ and $e_p=v_p(U)$,
\begin{equation}\label{eq:pell-lucas-exact-staircase}
                         v_p(D_r)=2re_p.
\end{equation}
At the last step,
\[
                         D_\theta=32^\theta U^{\ell-1}.
\]
\end{theorem}

\begin{proof}
Factoring $U^{2r}$ from the exact tail proves the first two statements in
\eqref{eq:pell-lucas-tail-factorization}.  The leading coefficient is
coprime to $U$ by Lemma~\ref{lem:pell-lucas-coefficient-unit}; adding a
multiple of $U^2$ preserves that coprimality.  This proves the exact
valuation.  The final formula uses $c_\theta=1$.
\end{proof}

For a hypothetical squarefull packet, every $e_p\ge2$, and at an
odd-kernel prime $e_p\ge3$.  Thus the $r$th deviation has exact depth at
least $4r$, respectively $6r$, at those primes.  This is an all-order
necessary condition, not a contradiction.

The first tail also couples to the companion.  Define
\[
 W=\frac{u_{\ell^2}/u_\ell-\ell}{U^2},
 \qquad
 S=\frac{v_{\ell^2}-v_\ell}{U^2}.
\]
The fourth-order formulas in
\cite[Corollary~5.2]{ZhangLucas2026} yield
\begin{align}
 W&\equiv\frac{4\ell(\ell^2-1)}3\pmod{U^2},
 \label{eq:pell-lucas-W}\\
 S&\equiv4v_\ell(\ell^2-1)\pmod{U^2}.
 \label{eq:pell-lucas-S}
\end{align}
In particular $W$ is a support unit.

\begin{theorem}[Paired companion splitter]
\label{thm:pell-lucas-companion-splitter}
One has
\begin{equation}\label{eq:pell-lucas-paired-correction}
                  \ell S\equiv3v_\ell W\pmod{U^2},
\end{equation}
and therefore
\begin{equation}\label{eq:pell-lucas-channel-split}
 \ell S\equiv6W\pmod{A^2},
 \qquad
 \ell S\equiv-6W\pmod{B^2}.
\end{equation}
The factor $6W$ is invertible modulo $U^2$.  The unique $Z\pmod{U^2}$
satisfying
\[
                         6WZ\equiv\ell S\pmod{U^2}
\]
obeys
\[
 Z\equiv A_{2\ell}\pmod{U^2},\qquad
 Z\equiv1\pmod{A^2},\qquad Z\equiv-1\pmod{B^2},
 \qquad Z^2\equiv1\pmod{U^2}.
\]
\end{theorem}

\begin{proof}
Multiply \eqref{eq:pell-lucas-S} by $\ell$ and
\eqref{eq:pell-lucas-W} by $3v_\ell$; their right sides coincide.  Next,
\[
 \frac{v_\ell}{2}=A_{2\ell}=2A^2+1=4B^2-1,
\]
which gives the opposite channel signs in
\eqref{eq:pell-lucas-channel-split}.  Every prime divisor of $U$ exceeds
$\ell+1$ and is therefore coprime to $6W$.  Cancellation shows that the
splitter residue is $A_{2\ell}$.  Since $A^2-2B^2=-1$, one has
$\gcd(A,B)=1$; hence its two squared signs combine modulo
$A^2B^2=U^2$ and give the last congruence.
\end{proof}

There is also a precise negative result.  A single normalized correction is
not rigid.  Consider the full-premise claim that every positive odd
$k\equiv1\pmod{u_3^2}$ satisfies
\[
              \frac{u_{3k}/u_3-1}{u_3^2}\equiv0\pmod5.
\]
It is false.  Since $u_3=35$, take
\[
                         k=2451=1+2\cdot35^2.
\]
Two independent recurrence computations, one linear and one by binary
powering in $\mathbf Z[T]/(T^2-6T+1)$, give
\[
 u_{7353}\equiv85785=35\cdot2451
       \pmod{214375=35(5\cdot35^2)}.
\]
After the justified divisions, the normalized correction is $2$ modulo
five.  This is an explicit instance of Zhang's local-surjectivity theorem
\cite[Theorem~5.6]{ZhangLucas2026}.  It retires only the exact fixed-zero
subclaim.  The all-order staircase and paired companion route remain
active because that theorem does not independently prescribe both
corrections.

The Lean companion kernel-checks the even-power tail algebra, coefficient
support argument, paired congruence, channel signs, modular cancellation,
square-root reconstruction, and the local-surjectivity logical boundary.
It also checks the displayed recurrence residue by kernel reduction, without
\texttt{native\_decide}.  The remaining sufficient target is a correlated
opposite-channel depth-three exclusion satisfying this complete staircase,
the earlier kernel character triangle, and the third-order factorization
ledger.  No full-premise packet counterexample was found, so that route is
not retired.
